\subsection{RSSI Measurement and Signal Conditioning}

The performance of the antenna alignment algorithm is very sensitive to
accurate signal strength readings. RSSI or Received Signal Strength Indicator is
the metric for measuring the radio strength of incoming packets in dBm. Higher RSSI
means higher signal strength which could be leveraged to measure the alignment.

The raw RSSI readings, however, have multipaths, fading, packet-to-packet
differences and other interference that makes them quite noisy to use directly. Thus,
a simple signal conditioning is applied to the input before applying it to the
algorithms and policies.

\subsubsection{RSSI Sampling}

RSSI data is collected from the integrated Wi-Fi radio on the ESP32-S3 board from
packets collected by the directional antenna. The RSSI measurements are performed
for each orientation (azimuth and tilt) in sequence. In the actual implementation
of the task, there were taken 20 RSSI measurements per angle with 10 ms interval in
between, which amounts to 200 ms per decision step.

\subsubsection{Noise Elimination and Filtering}

RSSI readings vary widely even if antenna orientation is kept constant. For reliable
control and learning, the following filters have been employed:
\begin{itemize}
    \item \textbf{Temporal averaging:} Averaging successive RSSI readings from the
    same orientation to eliminate noise
    \item \textbf{Dead-zone trend detection:} The difference between successive
    average RSSI values is used to classify whether there was any significant change
    in signal strength
\end{itemize}

These preprocessing techniques enable the agent to react to trends and not noise.

\subsubsection{Feature Extraction for State Space}

To efficiently capture the environment, an RL agent needs to build a compact state
representation as vectors. The state vector will be built using:
\begin{itemize}
    \item \textbf{Antenna Orientation:} Azimuth is discretized using increments of
    $10^{\circ}$ from $0^{\circ}$ to $180^{\circ}$. Tilt is discretized using increments
    of $5^{\circ}$ from $25^{\circ}$ to $75^{\circ}$.
    \item \textbf{Signal Trend:} Rate of change in averaged RSSI values is classified
    into three different categories:
        \begin{itemize}
            \item \textbf{Rising:} $\Delta RSSI > +0.8$ dB
            \item \textbf{Constant:} $-0.8 \leq \Delta RSSI \leq +0.8$ dB
            \item \textbf{Falling:} $\Delta RSSI < -0.8$ dB
        \end{itemize}
\end{itemize}

The resulting state vector will be represented by the following equation:
\begin{equation}
\mathbf{s} = \left( i_{\text{az}},\, i_{\text{tilt}},\, \Delta \mathrm{RSSI} \right)
\label{eq:state_vector}
\end{equation}
This approach combines both the antenna orientation and signal trend into the state
vector without increasing its size too much. Consequently, the Q-table can be deployed
on the ESP32-S3.

\subsubsection{Data Logging}

To facilitate debugging, performance evaluation, and further analysis, RSSI readings are
stored along with timestamp, azimuth and tilt index, chosen action, and trial ID.